\paragraph{Abstract}

Diese Hausarbeit beschreibt die Entwicklung einer grafischen Benutzeroberfläche (GUI) für Poker-Spiele mit PyQt6 als Ergänzung zur Bachelorarbeit über Counterfactual Regret Minimization (CFR).
Die GUI ermöglicht zwei Spielmodi: Agent vs Human, bei dem Spieler gegen trainierte KI-Strategien antreten, und Human vs Human, ein Multiplayer-Spiel über Netzwerk.

Die Implementierung demonstriert fortgeschrittene Python-Konzepte wie das Signal-Slot-Pattern von PyQt6, die Integration von asyncio mit QThread für WebSocket-Kommunikation, Thread-Safety-Mechanismen, umfassende Tests mit pytest sowie die Verwendung von Pickle und Kompression für die Serialisierung trainierter Strategien.

Die Architektur trennt klar zwischen GUI-Schicht, Client-Schicht, Server-Schicht und Logik-Schicht und ermöglicht die Wiederverwendung der Game-Logik aus der Bachelorarbeit.
Die Arbeit zeigt erfolgreich, wie moderne Python-Frameworks und -Patterns eingesetzt werden können, um eine funktionsfähige, wartbare und getestete Anwendung zu entwickeln.
